% ============================================================
%  Niels Treschow, "Gives der noget Begreb eller nogen Idee om
%  enslige Ting? Besvaret med Hensyn til Menneskeværd og Menneskevel."
%  In: Det Kongelige Danske Videnskabers Selskabs Skrifter, femte Deel (V),
%      1. Hæfte. Signed "1807"; volume published Kjøbenhavn 1810.
%  TRANSCRIPTION (Danish) — from Internet Archive scan `kongeligedanskev35kong`
%
%  Scan: ~/bibliotek/Treschow, Niels/videnskabernes-selskabs-skrifter-ser3-bd5-1810.pdf
%        (632 PDF pp.; ser.3 bd.5, 1810; MBLWHOI copy, 300 ppi, public domain).
%  TYPE: ANTIQUA (roman), NOT Fraktur — learned-society Skrifter convention.
%        Emphasis = antiqua ITALIC -> \emph{}. OCR with -l dan (or the ABBYY
%        layer as base), then image-verify every page.
%  ORTHOGRAPHY: this printing uses "ö" (o-diaeresis), NOT "ø" — preserve exactly
%        (Undersögelsen, Spörgsmaal, nödvendigt, Skiönhed, Udförelse). Also
%        gi- for gj-/gie- (Gienstande, giöre). Match the glyph.
%  OFFSET (VERIFIED): PDF = printed + 10.
%        Title page (printed 223) = PDF 233; body printed 225 = PDF 235;
%        last text page printed 254 = PDF 264 (ends "...gandske passende." + rule).
%  EXTENT: continuous essay, printed pp. 223--254 (~32 pp.), no chapter divisions.
%        PDF 265 onward = volume Sagregister (index), NOT part of the essay.
%  No Rettelser/errata leaf applies to this essay.
%  See RESUME-NOTES.md for the OCR pipeline, page map, and progress.
% ============================================================
\documentclass[12pt,a4paper]{article}
\usepackage[utf8]{inputenc}
\usepackage[T1]{fontenc}
\usepackage{libertinus}
\usepackage{libertinust1math}
\usepackage{textalpha}   % Greek words typed directly where they occur, if any
\usepackage[danish]{babel}
\usepackage[protrusion=true,expansion=true]{microtype}
\usepackage{setspace}
\usepackage[margin=1.2in]{geometry}
\usepackage{fancyhdr}
\usepackage{hyperref}

\hypersetup{
  pdftitle={Gives der noget Begreb eller nogen Idee om enslige Ting?},
  pdfauthor={Niels Treschow},
  hidelinks
}

\pagestyle{fancy}
\fancyhf{}
\rhead{Treschow, \emph{Om enslige Ting} (1807/1810)}
\cfoot{\thepage}

\setstretch{1.3}

% Original printed page numbers are recorded in comments "% printed p. NN (PDF NNN)"
% at each page boundary, matching the collation practice in the other transcriptions.

\title{\textbf{Gives der noget Begreb eller nogen Idee\\ om enslige Ting?}\\[6pt]
       \large Besvaret med Hensyn til\\[4pt]
       \large Menneskeværd og Menneskevel}
\author{Af Professor \textsc{Treschow}}
\date{\small Det Kongelige Danske Videnskabers Selskabs Skrifter,\\
      femte Deel, 1.~Hæfte. 1807 [udg.\ Kjøbenhavn 1810]}

\begin{document}
\maketitle
\thispagestyle{fancy}

\bigskip

% === TEXT BEGINS HERE (body: printed p. 225 = PDF 235) ===

% printed p. 225 (PDF 235)
Inden jeg begiver mig til Undersögelsen af det her opgivne Spörgsmaal
er det nödvendigt at forklare Meningen og tillige Vigtigheden deraf.
Ved Begreb forstaaer jeg med mange nyere Philosopher mere end den blotte
Forestilling. Denne kan være sandselig, hentet af umiddelbar Erfarenhed,
der kun omfatter Tingenes Overflade, ei det Indvortes og Bestandige deri.
Begrebernes Gienstande ere tvertimod Tingenes væsentlige og uforanderlige
Egenskaber, og selv enten, som Slægter og Arter, ved blot Abstraction af
Erfarenhed dannede eller have ved nogen anden Tænkekraftens Virksomhed
efter dens egne Love faaet en baade fast og eiendommelig Form. Denne Form
bestaaer deri, at Begreber tilligemed en vis Stof fremstille et Forhold
mellem det af Erfarenhed Bekiendte og noget Ubekiendt; Lighed f.~Ex.\ en
fælles Egenskab for flere i övrigt forskiellige Ting, Aarsag eller Kraft,
en indvortes Beskaffenhed, der indeholder Grunden
% printed p. 226 (PDF 236)
til det Udvortes og Synlige. Ideer igien forestille de höieste
Principier, som deels kan være det yderste Maal for vore Handlinger,
deels de almindeligste og förste Realgrunde til Gienstandene selv. Begge
ere uendelige, og derfor kun ved en uophörlig Nærmelse muelige at naae,
disse i vore theoretiske Grandskninger, i det vi forfölge Aarsagernes
lange Række til det överste Led, disse i vore Handlinger, som, hvor
ypperlige de end kan være, dog for Idealets Fuldkommenhed staae langt
tilbage.

Baade Begreber og Ideer have det tilfælles, at de indeholde noget
Bestandigt eller Væsentligt, istedenfor at Sandsernes Forestillinger kun
underrette os om de Forandringer, vi selv saavelsom alle udvortes Ting
ere underkastede. Men enslige Gienstande ere Erfarenheds, det er
Sandsernes egentlige Formaale: det Bestandige og Almindelige i samme
tilegner Forstanden eller Fornuften sig, og overlader Videnskaben at
ordne, undersöge, bearbeide: hvad der bliver tilbage pleier den, som et
unyttigt Skal, hvoraf Kiernen er udpillet, at forkaste.

Det enslige Menneske er ingen Undtagelse fra denne Regel. De
Videnskaber, der giöre vor egen Art til deres egentlige Formaal, saasom
Anthropologie, Moral, Politik og Lovgivning, betragte kun Menneskheden i
Almindelighed, ikke Mennesket. De mest practiske giöre det Hele, det
Almindelige til Öiemed, og tilsidesætte Delene, eller holde sig
berettigede til at ansee dem som blotte Midler for hiint. Derfor opoffres
saa mange Tusende for et saa kaldt almindeligt Bedste: paa ingen Altere
ryger Menneskeblod tiere end paa dem, som ere Menneskheden
% printed p. 227 (PDF 237)
selv indviede. Thi de enslige Ting, tænker man, ere forgængelige, Arten
eller Slægten allene forgaaer ikke. De Egenskaber, hvorved et Individ
udmærker sig fra et andet, ere blot empiriske: men have de ingen
uforanderlig og væsentlig Charakter, kan man derom ei giöre sig noget
fast Begreb; hvad ere de da vel andet end et Slags Luftsyner eller
Göglebilleder, behagelige nok maaskee for en kort Tid at anskue; men
hvorfor skaane dem, hvis de ere nogen Plan for Evigheden til Hinder,
efterdi de dog ere bestemte til i næste Öieblik at forsvinde?

Man har herimod paa den anden Side sagt, at Slægt og Art dog ei bestaae
uden ved enslige Ting, og at hine derfor tilligemed disse maae forgaae.
Dette Svar, der gandske stemmer overens med den af de fleste Philosopher
antagne Mening, er dog maaskee mindre tilfredsstillende end skikket til at
stoppe Munden paa dem, der uden nogen dybttænkt Theorie vil forsvare den
Uret, de under Skin af det Heles Vel tilföie saa mange blandt deres
Medmennesker. At udgive Arten for et blot Begreb uden besynderlig
Gienstand vilde være den letteste Maade at afgiöre Sagen paa, dersom den
gandske passede til alle de Særsyner, man derved vil oplyse. Dertil kan
for det förste regnes den paafaldende Lighed, man ei kan undgaae at
bemærke mellem Planter og Dyr af samme Art; dernæst de Kræfter, der hos
alle eller mange Slags Væsener altid virke efter de samme Love. Thi
hverken ville Naturbeskriverne tilstaae, at deres naturlige Arter ei have
nogen anden Grund end Sindets Abstraction: ei heller vil man under den
Forudsætning, at Naturkræfterne ere blot almindelige Begreber, kunne
forstaae hvorledes
% printed p. 228 (PDF 238)
samme Kraft, f.~Ex.\ den til Attraction og Repulsion, efter de adskillige
Forestillingsarter, man derom har udtænkt, enten ved et Slags Deling paa
engang kan findes hos mange enslige Ting, eller hvorledes der af en Kraft,
som kun i Henseende til Arten er den samme, kan være ligesaa mange
individuelle, som Substanser, der besidde den. Det sidste forekom
\emph{Aristoteles}, blandt flere saadanne Exempler, som det synes, saa
utænkeligt, at, endskiöndt han ellers om Begrebernes Oprindelse forkastede
Platos Mening og var enig med de för anförte empiriske Philosopher, saa
antager han dog en eneste almindelig virksom Forstand, der hos hvert
Menneske vel yttrer sig paa en besynderlig Maade, men numerisk ei er
forskiellig. Man begriber kun ikke, hvorledes han ei indsaae, at det
samme om enhver anden Kraft, der ligger enten i den menneskelige eller
fælles Natur, da ligesaavel maae gielde.

Det kan fölgelig ei ansees for en saa afgiort Sag, at Forstandens
Begreber, f.~Ex.\ det om Menneskheden i det Hele, uden for Individerne
selv, ingen virkelige Gienstande kan have. De platoniske Ideer, saavidt
de skal være Former eller Mönstre, hvorefter alle, der höre til en Art,
ligesom Aftryk ere dannede, har hverken Aristoteles, Locke eller Kant
tilstrækkelig giendrevet, ja, ei engang udtrykkelig modsagt, endskiöndt de
tydelig nok have viist Urigtigheden i mange tildeels falske eller skæve
Udtydninger af denne Theorie. I det practiske Liv, i den almindelige
Menneskeforstands egne Domme, i Lovgiveres og Statsmænds Planer viser sig
endog da denne Menings Indflydelse, naar de, der antage eller forudsætte
den, neppe ere sig dette
% printed p. 229 (PDF 239)
selv bevidste, eller derom have giort sig noget tydeligt Begreb. Spörger
man efter disse Ideers Gienstande, saa synes det ei heller urimeligt, at
de kan ligge baade i os selv og den hele udvortes Natur, ligesom under en
Taage eller et blindende Skin, begravne, og at der til ikke blot at ahne
dem, som de fleste giöre, men til tydelig at blive dem vaer udfordres et
skarpere Syn, som, enten det er givet af Naturen eller erhvervet ved egen
Flid, dog kun findes hos nogle faa.

Men om endog denne Theorie, saavidt den strækker sig, er rigtig, kunde
den dog ei maaskee være alt for indskrænket og falsk, ifald den derhos er
udelukkende? Gives der ei tillige saavel Ideer som Begreber om enslige
Ting, eller ere de om Arter og Slægter alene muelige? Dette Spörgsmaal er
fornemmelig, naar det anvendes paa Mennesker, overmaade vigtigt. Det vil
nemlig sige saa meget: Har ethvert Menneske, ligesaavel som Arten selv, en
vis Grundform eller væsentlig Character, der under alle Livets
Omstændigheder bliver den samme? Er Individualiteten noget andet end et
Spil af vexlende Skikkelser eller nye Bestemmelser, hvori intet er
uforgængeligt, undtagen hvad enhver har tilfælleds med alle andre, det er
med den almindelige Menneskenatur? eller ere disse Forandringer kun
successive Udviklinger af en Kierne, der, hvad enten den först begynder at
groe eller allerede er bleven til et frugtbringende Træe, eller efter
flere Forvandlinger er bleven endog sig selv ukiendelig, dog stedse
beholder de oprindelige Grundtræk tilbage? Er endelig Maalet, hvortil
enhver endog ubevidst sigter, det Mönster, hvorefter han maae dannes, i
Sandhed en uendelig Idee, som uagtet al Mangfoldighed,
% printed p. 230 (PDF 240)
ja, tilsyneladende Strid har indvortes Enhed og urokkelig Fasthed i sig
selv? Det er klart, at kun i dette Tilfælde har ethvert Individ et
ubeskriveligt Værd, i det andet derimod Arten allene. At denne Forskiel
paa Meninger nödvendig maae have særdeles vigtige practiske Fölger, er
næsten overflödigt at erindre.

Efter den almindelige Metaphysiks Fremstilling bestaaer Enslighedens
Kiendemærke i en fuldkommen Bestemthed saavel i Henseende til enhver
Beskaffenhed, man kan tillægge Tingene, som i Henseende til Tid og Sted.
Da nu begge idelig forandres, saa synes deraf at fölge, at hos enslige
Væsener, som saadanne er intet uforanderligt. Men dette Kiendemærke er
ogsaa blot logisk, og vil ei sige andet end at man, for at skille saadanne
Ting fra andre af samme Art, ei har andet Middel end Beskrivelser, fordi
deres logiske Væsen selv, som Definitioner kan bestemme, end sige deres
metaphysiske og reale, er os aldeles ubekiendt, da vi dog hvad almindelige
Begreber angaaer i det mindste kiende det förste.

Skolastikerne komme Hovedsagen meget nærmere, og deres Dispyter om enslige
Tings Identitet ligger i Materiens eller Formens Bestandighed, röbe baade
Skarpsind og Indsigt. Den förste er fælles for mange: skulde Stoffens
Eenhed være nok og Formen ei tages i nogen Betragtning; saa kunde
Begrebet, hvis Væsen er Formen, umuelig blive det samme. Thi af samme
Stykke Træe kan hugges en Klods, og Billedet af en Gud; Formen er derimod
i det Hele foranderlig, med mindre der gives visse Love, som bestemme den
for ethvert Individ og indskrænke dens Foranderlighed til saadanne
Grændser, at den indvortes dannende Kraft efter de Ensliges besynderlige
Natur eller
% printed p. 231 (PDF 241)
Idee i en bestandig Fremgang til större og större Fuldkommenhed kan
udvikle sig, men at disse derimod ifölge en organisk Forbindelse mellem
alle dertil vexelviis vække og hielpe hverandre.

Efter Platos Mening vare Ideerne, ligesom Mönstre, hvoraf mange Kopier
blive tagne, almindelige Begreber; og de enslige Tings Forskiel beroede
ligeledes ei heller paa Ideens, men paa Materiens större eller mindre
Beqvemhed til at udtrykke den. De nyere Platonikere vare af samme Mening.
Det Schellingske System, som i Grunden ei er andet end det nyeplatoniske i
en efter Naturvidenskabernes nærværende Tilstand forandret Skikkelse,
sætter ogsaa en Ære i at hæve sig over enslige som sandselige Gienstande
eller vaklende Billeder af det Absolute og af Ideerne eller dets evige og
uforanderlige Former. Individuerne bör derfor selv stedse stræbe at vende
tilbage til den almindelige Natur: og skiöndt deraf bestandig nye
fremtræde paa Skuepladsen, ere de dog ikke mere de samme. De besynderlige
Siele ere Smaadele af det store Ocean, som först forvandlede til Dunster
eller Taage, Legemverdenens Sindbillede, dernæst til Regndraaber, samle
sig i Kilder, Bække og Floder, for efter en kort afsondret Tilværelse,
hvis Nydelser Savn og Længsel dog idelig forbittre, at fortryde sit Affald
og igien blive salig ved Forglemmelse.

Er denne Forestilling rigtig, saa har det menneskelige Liv ei alene liden
eller ingen Værd, men er endog et virkeligt Onde, fra hvilket man maae
stræbe jo för jo heller at blive frie. Legemet, Sandseligheden og alt hvad
der hörer til vor Individualitet er et forgængeligt Klædebond, der ikkun
hemmer Sindets frie
% printed p. 232 (PDF 242)
Bevægelser, en Vægt der nedtrykker os til Jorden, eller et Fængsel,
giennem hvis tykke Mure Lysets Gienskin fra den aabne Natur ei kan
indtrænge. Vel have vi ikke Lov til at brække ud, vi maae blive indtil vi
enten blive frikiendte eller aflöste: men hvor meget mere velgiörende er
dog ei under denne Forudsætning Grusomheds og Blodgierrigheds Stemme, som,
i det den förer Tusende til Slagterbænken, tillige aabner dem den sande
Frieheds Porte end Mildheds og Menneskeligheds, der ved at raabe Tilgivelse
egentlig fængsler Synderen paa nye? Kunde man vel önske, at denne Theorie
vandt almindeligt Bifald, og at Regenter og Lovgivere derover maatte troe
sig berettigede til at behandle Mennesker som blotte Midler, bruge dem som
Materialier til et eller andet Pragtverk, som de agtede at opföre? Er ei
Begrebet om Ret og Uret aldeles grundet paa Personlighed, d.~e.\
Individualitet? Ophæv denne, sæt at denne kunde forstaaes uden noget
Selvstændigt og Bestandigt, at Personer kun vare flygtige Skygger, der
komme og forgaae, hvorledes kunde nogen Rettighed da tilkomme dem, som jo
i næste Öieblik maatte være forloren? hvorledes kunde man straffe eller
belönne dem, hvis de selv forsvandt tilligemed den Handling, man tilregnede
dem? Alle naturlige saavelsom borgerlige Love forudsætte noget Bestandigt
hos Personer: de vilde være ubillige og ufornuftige, dersom Forholdet
mellem Enslige ei var ligesaa varigt som bestemt.

Spörger man i denne Sag alene Erfarenhed tilraads, faaer man, som i de
fleste Tilfælde, hvor det kommer an paa almindelig Gyldighed, et gandske
tvetydigt Svar. I den hele uorganiske Natur findes intet af de der
saakaldte enslige Væsener, der
% printed p. 233 (PDF 243)
besidder nogen uforanderlig Charakter. Arterne selv ere saa vaklende, at
faste Begreber derom ere meget vanskelige at erholde. Ethvert uorganisk
Legeme er et Aggregat af mange, enhver Deel kan bestaae for sig selv
ligesaavel som i det Hele uden at tabe sin eiendommelige Natur. Ved Deling
kan tillige Formen forandres, der ligesaa lidet er dem væsentlig som
Materien: hvem kan derfor sige, hvori deres Individualitet bestaaer? Naar
ophörer Floden eller Bygningen vel at være den samme? Exemplet af
Theseus's Skib er bekiendt og med Anvendelse paa denne Sag tidt blevet
anfört. Vare Krystallernes Grundformer, som i denne Deel af Naturen dog
give de bedste Kiendemærker, end uforanderlige, saa ere de dog, ligesom
Anaxagoras's Homoiomerier delelige, maaskee i det uendelige; altsaa
ligeledes Aggregater, ei sande enslige Væsener. De virkelige Individuer
finde vi kun i den organiske Natur. Jo fuldkomnere Organisationen er,
desto större Bestandighed er der i deres Form, desto mindre kan nogen Deel
savnes, nogen Omsætning skee uden væsentlig Forandring i det Hele.

Blandt alle Figurer er Cirkelen den fuldkomneste, fordi enhver Punkt deri
er gandske bestemt, og alle formedelst Middelpunkten paa samme Maade.
Ingen Forskiel mellem enslige Cirkler uden i Störrelse. Tænke vi os denne
ubegrændset, hvilket ei gaaer an i nogen anden Figur, da Linierne fra
Middelpunkten til Enderne i disse altid ere af ulige Störrelse, og fölgelig
ei uendelige; saa kan der af Arten selv ei være flere end en eneste. Alle
andre Slags krumme og retliniede Figurer kan ogsaa i andre Henseender,
f.~Ex.\ i Ordinaternes, Sidernes eller Vink%
% printed p. 234 (PDF 244)
lernes Forhold til hverandre, være forskiellige. Ei uden Aarsag
fremstillede derfor de Gamle det höieste Realprincip, Guddommen som den
fuldkomneste Individualitet, under Billedet af en Cirkel eller Kugle.

Naturens hele Stræben, den almindelige Fuldkommenheds-Drivt har intet
andet Formaal en Regelmæssighed og med samme tillige Bestandighed i
% sic: "en" printed for "end" (dropped d) — translated as "than"
Formen. Moralsk saavelsom physisk Fuldkommenhed bestaaer i uforstyrret
Harmonie eller, med andre Ord, fuldstændig Individualitet, hvorved alle
Dele ere saaledes forbundne, at ingen er mindre væsentlig eller nödvendig
end den anden, og at ingen Tilfældighed deri har Sted.

Drivten til Vedligeholdelse er af Naturen mest umiddelbar henrettet paa
det enslige, middelbar allene paa ethvert Væsens egen Art. Den oprindelige
Lyst til Virksomhed er vel for saavidt ingenlunde blot egennyttig, som
Stoffen, den nöder os til at bearbeide og efter en Idee at danne, ei er
vor egen Natur alene og hvad dertil henhörer eller hvad der tiener til
sammes Fuldkommenhed, men andre Menneskers tillige, ja, den hele
Udenverdens, naar den kun er skikket til at modtage saadan Dannelse, kan
være dens Formaal. Vi elske Skiönhed, ei fordi den findes i vor egen
Person, men ogsaa naar vi i hvilkensomhelst anden Ting blive den vaer,
allevegne söge vi derfor at realisere dens Idee. Ikke de mindre er Grunden
% sic: "Ikke de mindre" printed for "Ikke desmindre" — translated "Nevertheless"
til det behagelige Indtryk den giör, at vi derved komme til Bevidsthed
eller Selvfölelse, d.~e.\ Fornemmelse af vor egen Kraft, dens Tilvæxt eller
Styrke.

% printed p. 235 (PDF 245)
Selvbevisthed er den fornemste Yttring af Individualitet, det höieste os
bekiendte Trin, som noget endeligt Væsen har besteget. Jo klarere og
tydeligere vi adskille os fra andre, kiende hvad der udmærker os fra dem,
desto fuldkomnere ere vi. Al Philosophie bestaaer i Kundskab om os selv,
ei efter den fælleds Menneskenatur allene, men især den individuelle.
Frugten af Lærdom og Eftertanke er fornemmelig en Art af Besindelse. Al
Udvikling og Oplösning af Begreber, al Sammensætning og Forbindelse af de
samme tiener kun til at giöre Bevidstheden gandske klar, sætte os istand
til at giöre os fuldkommen Rede for os selv. Samvittighed og Sædelighed
beroe ei heller paa noget andet. At handle ei af Drivt, men af Principier,
er det ei at handle med tydelig Bevidsthed af samme? Thi ogsaa hos dem, der
bestemmes af Fölelser, ere Grundsætninger virksomme, skiöndt efter dunkle
og utydelige Forestillinger. Vort Væsens hele Fortræffelighed bestaaer
altsaa i Individualitet, Formaalet for al Stræben er at opnaae samme i den
allerhöieste Grad: den er fölgelig for det förste lovlig, det er hverken
nödvendigt, nyttigt eller mueligt at undertrykke den: dernæst er den det
Mönster, hvorefter enhver maae danne sig, den Idee, man stedse bör söge at
komme nærmere: hvoraf flyder, at Ideer kan være individuelle, og at deres
Gienstande, ligesaavel som Arterne, ere evige, uendelige og uforanderlige.

Man kan endog opkaste det Spörgsmaal, om der vel gives andre Ideer af
virkelige Ting end de anförte, og om den almindelige Mening, at disse
enten nödvendig maae være almindelige Begreber, eller i det mindste
ligesaavel kan have Arter
% printed p. 236 (PDF 246)
og Slægter til Gienstande som enslige Ting, ei grunder sig paa en
Forvirring af Tingenes logiske og ideale Væsen med det reale, som kun
Individuer virkelig kan besidde. Den Dispyt om Slægter og Arter uden for
Tankerne i Sandhed ere til, er, som bekiendt, meget gammel. De, der ansee
dem for virkelige, beraabe sig paa de faste Grændser, som Naturen synes at
have sat mellem Væsener af forskiellig Art, der giöre at ingen Overgang er
muelig fra den ene til den anden. Hvad de fuldkomnere Organismer saavel i
Plante- som Dyreriget angaaer, holdes Sagen af de fleste Naturkyndige for
afgiort, især er det menneskelige Kiön efter deres Tanker et talende Beviis
derpaa, da endog de ringeste blandt vor Art saa tydelig udmærker sig fra de
ypperligste blandt Naboarterne, eller, om man vil, Slægterne, at ingen
Forvirring eller Blanding uden af grov Uvidenhed her synes muelig. Ikke
desmindre gives der Kiendsgierninger og andre Grunde, der kunde opvække
Tvivl derimod. Jeg vil ei engang anföre mange usikkre, som af adskillige
berömte Mænd her ere brugte, ei Formodninger og Hypotheser, der af
anseelige Forfattere, saasom Monboddo, Hallen, Lamark, selv af vor
Fabricius, ere fremsatte, men ei tilstrækkelig beviste. For det förste er
det dog vist, at jo almindeligere Begreberne blive, jo mere de fierne sig
fra de individuelle, desto mere löbe deres Grændser sammen, desto
vanskeligere blive de at bestemme. Slægternes Kiendemærker ere i
Almindelighed mere vaklende end Arternes, og blandt hine mest de höiere,
nemlig Ordnernes, Familiernes og Classernes af organiske saavelsom
uorganiske Legemer. I det vor Kundskab om Naturen, endog den blot
empiriske, udvider sig, see vi efterhaanden de Grændser at forsvinde,
% printed p. 237 (PDF 247)
som vi troede ei vilkaarlig, men efter dens egen Anviisning at have
bestemt. For det andet naar end ikke den længste Tid selv synes at kunne
forrykke dem, saa er det alligevel tvivlsomt, om denne Bestandighed er
absolut og væsentlig, eller blot relativ, da vi nödes til at ansee for
urokkeligt hvad nogle faa Aartusender, til hvilke vor Erfarenhed, naar den
kommer til det höieste, dog maae indskrænke sig, har ladet staae. Vi finde
for det tredie dog ligesaavel Afarter, f.~Ex.\ Negre eller blandt Hundene
Mopse og Mynner, som egentlige Arter, der uagtet de formodentlig ved
tilfældige Omstændigheder ere blevne saa afvigende, tilsidst have antaget
en uudslettelig Charakter. Hvorledes vil man nu vel godtgiöre, at det
samme om de sidste ei ligeledes kan gielde?

Det er endelig ei saa vanskeligt at begribe, hvorledes Arter formedelst
den almindelige Tiltræknings og Affinitets Love, der i Plante- og
Dyreriget faaer adskillige nye Bestemmelser og en höiere Betydning, kan
være blevne til, om endog fra Begyndelsen lutter Individuer i mange,
skiöndt ei endnu saa tydelig bestemte Gradationer og Former af samme Kraft
have været frembragte. Thi da Lige overalt söge Lige, da Forbindelser,
Selskaber, Avledrivten, og Drivten til Efterlignelse stedse maae bringe
flere hverandre endnu nærmere end de alt af Naturen ere; saa maae alle
Enslige derover omsider have dannet Grupper eller Hobe, der vel paa
Grændserne kan staae i et dobbelt og tvetydigt Forvantskab, nemlig baade
til dem, der ere inden og dem, der ere uden for samme, men hvoraf dog den
störste Mængde kiendelig nok maae være adskilt fra dem, der höre til de om%
% printed p. 238 (PDF 248)
kringliggende. Oprindelig gives derfor kun Individuer; Arter og Slægter
synes derimod, ligesom Selskaber eller Kunstverker, at være af en anden og
senere Skabelse.

\emph{Linneus}, der ved et Slags Genieblik opfandt meget, som hos andre
kunde have været en Frugt af langvarig Eftertanke, meente at der i
Begyndelsen kun har været Slægter til, af hvis Forbindelse Arterne siden
ere fremkomne. Kants Hypothese om Oprindelsen til de forskiellige
Menneskeracer, hvortil Kiernen skal have ligget i det förste Menneskepar,
har rimeligviis hos ham havt samme Grund, skiöndt ingen af dem tydelig har
udviklet den. Jeg skulde troe, at begge have villet eller i det mindste
burdet derved sige omtrent fölgende: Hverken Slægter eller Arter have i
Begyndelsen været saa bestemte, som de ved Blandinger og Forbindelser
siden ere blevne. Formerne have, ligesom i Billedhuggernes Verker, for det
förste været gandske raae. Naturen har, saa at sige, ladet sig nöie med at
trække de store Grundlinier, efterhaanden ere de finere saavelsom den
sidste Politur kommen til. Det lader sig f.~Ex.\ tænke, at de förste Dyr
have været til i en Tilstand, hvori det kan have været vanskeligt at
bestemme, om Formen skulde forestille Mennesker eller et ufornuftigt Dyr.
Efterat hin Skikkelse begyndte at aabenbare sig, kunde det endnu være
uvist, hvilken Person især det var Mesterens Hensigt at afbilde. Finder man
dog overalt Physionomier, hvori det charakteristiske og individuelle synes
at mangle. Men jo mere Verket nærmede sig Fuldendelse, desto kiendeligere
bleve de udmærkende Træk. Naturen gaaer heri ligeledes kun gradvis frem.
Men det er klart, at Individualitet er dens
% printed p. 239 (PDF 249)
yderste Maal. Hvad det menneskelige Kiön angaaer, saa lærer Erfarenhed, at
Kulturen selv dertil bidrager, og at man af den Aarsag hos civiliserte
Folkeslags finder en Mangfoldighed og tillige Bestemthed i enslige
Personers Form, som blandt Barbarer og Vilde aldrig er saa stor. Det maae
fölgelig være Naturens Hensigt at individualisere sig; det Besynderlige
eller Individuelle er ikke Middel og Redskab for det Almindelige, men
Öiemedet i sig selv.

Men, om det endog var afgiort, at Individuerne ere Ideernes virkelige, ja
maaskee eneste Gienstande; har denne Sætning da vel nogen vigtig Fölge?
Faaer man derved nærmere Kundskab om Individuerne eller Tingene selv? Kan
man derom nu bedre end ellers giöre sig et tydeligt Begreb og saaledes
bestemme det, at vi blive satte istand til at adskille det ene fra det
andet? Hvis ikke, hvortil kan det Resultat nytte, som vi saa möisommelig
have sögt? Thi ligesaa lette som enslige Tings empiriske Kiendemærker maae
være at udfinde, som til Brug i det daglige Liv ere tilstrækkelige, saa
vanskelig er Grundformen selv enten i Sindet at fastsætte eller med Ord at
udtrykke. Ikke desmindre svæver os deraf et klart skiöndt forvirret Billede
for Öinene: vi vilde ellers af Omgang og Anskuelse ei bedre lære Tingene at
kiende end af blotte Beskrivelser. Malernes og Billedhuggernes Kunst
bestaaer i at fatte disse Grundformer og saaledes at fremstille dem, at de
endog ved Forskiönnelse snarere blive satte i större Lys end fordunklede.
Saaledes er det mueligt af ethvert menneskeligt Ansigt, af enhver ei
gandske fuldkommen Charakter at danne baade et Ideal og en Carricatur uden
at udslette eller giöre Grundtrækene mindre kiendelige paa den ene end paa
den anden
% printed p. 240 (PDF 250)
Maade. Hverken Ungdommens Roser eller Alderdommens Rynker forandre denne
Form: ere vi kun ret opmærksomme, blive vi den i begge vaer. Alder og
Kultur udvikle den mere og mere: i den tidlige Barndom er den gandske
ubestemt; Raahed og Vildhed giöre den ligeledes vanskelig at bemærke: i
hvilke Tilstande Individuerne ogsaa synes mest at ligne saavel dem, der
regnes til samme Art, som dem, der höre til forskiellige, f.~Ex.\ Mennesker
og Aber.

Det giör intet til Hovedsagen, at vi ei ere istand til hverken i
Beskrivelser eller Definitioner at fremstille denne Form. De förste tages
af Öieblikkets Fornemmelser, eller af de Indtryk, Tingene giöre paa vore
Sandser; hvori det Væsentlige og Tilfældige, det Objective og Subjective
stedse ere blandede. Til Definitioner mangler Sproget den fornödne Rigdom
af Ord, da det af Nödvendighed maae indskrænke sig til at betegne
almindelige Begreber: og endskiöndt det ogsaa har Navne paa enslige
Gienstande; saa har det dog ingen paa de individuelle Beskaffenheder, der
giöre dem til saadanne. Den menneskelige Forstands Indskrænkning maae nöies
med at kiende Tingene overhovedet og i det Store: men man bör derfor ikke
glemme, at de almindelige Begreber kun ere Hielpemidler til Oversyn, ei til
nöiagtig Kundskab om Tingene i sig selv. Thi i den virkelige Verden findes
kun Individuer, og deres Væsen er ikke det, der udtrykkes ved Artens eller
Slægtens Kiendemærker, men hvad der ligger i den evige Idee, hvoraf utallige
Modificationer, der i Tiden vise sig, ere muelige. Denne Idee er ei i den
Henseende almindelig, at den omfatter flere Ting; men for saavidt som den
udtrykker alle muelige Tilstande af den samme, kunde man maaskee give den
dette Navn: rigtigen er
% printed p. 241 (PDF 251)
det alligevel at giöre Forskiel paa denne abstracte Forestilling om det
Enslige og den concrete.

Min Hensigt er det ikke at kalde de Grundsætninger i Tvivl, paa hvilke
Naturbeskrivelsen bygger. Jeg erkiender med alle andre deres Nödvendighed
og Nytte. Men man vilde af Kiærlighed til denne Videnskab efter mine Tanker
gaae for vidt, hvis man enten blot af den Aarsag ansaae Arternes
Kiendemærker for væsentlige og uforanderlige, fordi den indskrænkede
Erfarenhed, vi have, ei udtrykkelig vidner derimod, eller om man i den
almindelige Betragtning vilde skiule det uvisse, vaklende, de mange
Undtagelser og Afvigelser, der ogsaa i Henseender til Arterne virkekelig
har Sted. % sic: "virke-|kelig" (linjedelt dittografi) for "virkelig"
Var alligevel Meningen om deres absolute Uforanderlighed end fuldkommen
beviist; saa fulgte deraf dog ikke, at Individuernes indvortes Charakter er
mindre bestandig; tvertimod tillader Analogien, som ovenfor er bemærket, os
at slutte saaledes: Ligesom Arternes Kiendemærker ere mindre vaklende end
Slægternes og Familiernes, saa ere Individuerne, som de laveste Begreber,
endnu mindre vigtige Afvigelser fra deres Grundform underkastede. Den
Sætning, jeg her forsvarer, kan altsaa med Naturbeskrivernes Theorie meget
vel bestaae, endskiöndt denne Videnskab deraf ei kan betiene sig, fordi
hiin baade heri er umuelig at anvende, og, hvis det kunde skee, vilde före
til en Vidtlöftighed, der langt overgik de Grændser, som en menneskelig
Videnskab maae have. I Historien, der nödvendig maae befatte sig med
enslige Personer, finder derimod den her afhandlede Sætning saavel en
nærmere Anvendelse som Stadfæstelse, begge Dele i de Charaktertegninger,
som den ei uden et dybt
% printed p. 242 (PDF 252)
forskende Blik baade i den menneskelige Natur overalt og i dens mangehaande
finere og individuelle, men med Evighedens Præg udmærkede, Former er istand
til at udkaste.

Forgieves söger jeg alligevel ved disse og flere saadanne Grunde at
overtyde dem om min Menings Rigtighed, som troe, at Svælget mellem Os og
det Absolute, hvilket man dog bör stræbe at udfylde, derved bliver endnu
videre og dybere; med mindre det lykkes mig at vise, at denne Sætning ei
allene med deres Identitets System kan bestaae, men at dette endog
nödvendig förer til samme Resultat. Til den Ende skal jeg nu saa kort og
tydelig, som mueligt, forklare hvorledes Individuernes Forhold til det Ene
og Absolute saavelsom til de evige Ideer efter disse Philosophers egne
Grundsætninger bedst og lettest synes at lade sig begribe.

Paa hvor mange Maader man end kan fremstille og virkelig har fremstillet
denne Lærebygning, forekomme mig dog alle at stemme overeens i fölgende
Lærdomme. Det överste Realprincip kan ikke tænkes uden som reen og
uindskrænket Kraft, Virksomhed eller Handlen. Men, da det uden for sig ei
finder nogen Stof at virke paa, og en Skabelse af intet enten er modsigende
eller dog aldeles ubegribelig; saa kan denne Virksomhed ei sammenlignes med
nogen anden os bekiendt undtagen med Forstandens eller Fornuftens, saavidt
denne af sig selv og a priori frembringer Begreber og Ideer. Til disse kan
nu Stoffen heller ikke uden fra være given; den dannende Forstand maae
tillige være den frembringende Kraft; men dette kan den ikkun være ved
Selvansku%
% printed p. 243 (PDF 253)
else, eller med andre Ord: Stoffen til Ideerne kan den ikke tage af noget
andet end sig selv. Sandselighed kan man derfor ikke tillægge det
allerhöieste Væsen, fordi den altid maae betragtes som lidende og en blot
Evne til at blive rört. Derfor lære disse Philosopher, at der gives en
Forstandens Anskuelse, hvorom jeg paa dette Sted ei kan handle. Men heraf
fölger nu, at Stoffen til alle Ideer er den samme, al Forskiel imellem dem
beroer paa deres Form, d.~e., de mange Bestemmelser, eller egentlig,
Grændser, Indskrænkninger, hiin kan modtage. Enhver Idee er alligevel,
ligesom Cirkelen og alle andre Linier eller Figurer, uendelig i en vis
Henseende, fordi nemlig Störrelsen uagtet den bestemte Form er ubestemt. De
ere ligeledes evige og nödvendige, det er Tingenes eget Væsen de udtrykke.
Kun naar de i Tiden og Rummet, som Sandseverdenens Former, fremstille sig,
iförer deres evige Natur sig, saa at sige, et forgængeligt Klædebond; fra
deres himmelske Boeliger stige de ned i disse sublunariske Egne, hvor alle
Ting deels ere foranderlige deels vise sig paa deres Tilværelses laveste
Trin, og Ideerne i den meest fornedrede Skikkelse. Den iboende Kraft synes
enten, som i den grovere Materie, gandske död, eller, befinder sig, som i
Spirer, i sin Udviklings förste Begyndelse. Dette förer Tiden med sig, som
hverken kan tænkes uden Forandringer eller uden ordentlig Fölge, der ved en
Regel maae være fastsat. Men denne Fölge begynder med en Punkt, der er det
allermindste Reale og Tænkelige, voxer, og bliver omsider til en uendelig
Linie, eller saavidt man kan forestille sig Punkten som straalende, til
flere ja utallige. Paa samme Maade er enhver Kraft i Begyndelsen svag og
usynlig. Livet grændser til Döden förend
% printed p. 244 (PDF 254)
det med Tiden samler Styrke. Lys og Varme, de mest virksomme Stoffer,
skiule sig, som ulmende Gnister förend de blive vakte, under en dorsk
Materies Skikkelse. Men mod Formens Enhed, Ideens Uforanderlighed er saadan
successiv Fremgang ingenlunde stridig; tvertimod udgiöre Reglerne,
hvorefter den udvikler sig, en Deel af dens Væsen selv. Fremdeles: Saa
mange forskiellige Former, som der ere muelige, saa mange evige Ideer maae
der ogsaa være; thi Forstanden anskuer alt, dens Almagt kan giöre dem alle
virkelige; dens Godhed og Viisdom vil sikkert söge at bringe de almindelige
saakaldte virkelige Ting, d.~e.\ de i Tiden og Rummet tilværende, som dog
egentlig kun ere svage Billeder af de sande, nemlig Ideernes Gienstande,
disse Mönstre stedse nærmere.

Denne hele Theorie kan umuelig uden paa Individuer have nogen Anvendelse.
Thi for det förste er Arten ikke det Mönster, hvorefter disse maae danne
sig. Hverken er den saa bestemt, at nogen fast Regel deraf lader sig udlede,
ei heller hörer engang Fuldkommenhed til dens Begreb; da ethvert ensligt
Væsen efter den almindelige Mening, der giör Arten uforanderlig, hvor meget
det end afviger fra Regelen, dog nödvendig maae blive inden for Artens
Grændser. Dernæst ere ei Arterne alene, men ogsaa de Enslige i Henseende
til deres Form forskiellige. Det er derfor langt fra, at de förste udtömme
alle muelige Former; hvilket dog formedelst den uendelige Forstands Rigdom
eller Frugtbarhed nödvendig maae skee.

Jeg har uden at binde mig til nogen Forfatters eller Skoles Terminologie
sögt at fremstille Principierne til denne Lærebygning
% printed p. 245 (PDF 255)
paa saadan Maade, der forekommer mig paa engang at være den klareste og
tillige mest conseqvente. At mange baade saakaldte Platonikere og andre,
der bekiende sig til den samme, i at fordömme den naturlige Stræben, der
har den individuelle Fuldkommenhed eller det enslige Væsens Vedligeholdelse
til Formaal, have miskiendt Aanden af deres eget System, dertil er Aarsagen
deels denne, at de have anseet Individualitetens Kiendemærker allesammen for
ubestandige, uden at agte paa det Væsentlige og Uforanderlige, der
indeholder disses indvortes Grund; og til denne Vildfarelse ere de igien
forledede, fordi de ei have bemærket, at virkelige Individuer allene findes
i den organiske Natur, men dermed blandet de Aggregater, der udgiöre den
uorganiske. Deels have de, for efter deres Mening desto bedre at forklare
det Ondes Oprindelse, betragtet Ideernes Udgang, d.~e.\ Afbildning eller
Udvikling i Tiden og Rummet som et Slags Syndefald, der endogsaa hos de
nyeste Philosopher har været Anledning til selsomme Forestillinger;
endskiöndt dog nogle blandt dem, ja \emph{Plotin} selv, have indseet, at en
uendelig og bestandig Virksomhed ei er muelig uden ved endelige Producter
eller Skabninger, at disses Liv ei kan være andet end en uophörlig Stræben
efter det Bedre, og at de fölgelig alle maae være mere eller mindre onde,
d.~e.\ ufuldkomne.

Dog dette Slags Philosopher er det maaskee mindre vanskeligt herom at
overbevise end andre. Nogle iblandt dem komme os i denne Henseende selv i
Möde. \emph{Spinoza}, som jeg uagtet hans meget forskiellige
Fremstillingsart ei tager i Betænkning at sætte paa samme Liste, antager
f.~Ex.\ at enhver Menneskesiel
% printed p. 246 (PDF 256)
tænkt under Evighedens Idee er uforgængelig: hvormed han intet andet har
villet sige end at dens Væsen eller Grundform giennem alle Evigheder er den
samme, ei at den blot beholder Artens Kiendemærker. \emph{Fichte}, hvis
Lærebygning, som hans senere Skrivter udvise, stedse mere nærmer sig hiin
alexandrinsk og nu germanisk platoniske, er ei heller af den Mening, at man
ved at aflægge sin Individualitet vilde hæve sig til en större eller den
engang ulykkeligviis forlorne Fuldkommenhed. Ja, endog efter den seneste
Naturphilosophies Grundsætninger er Stræben at individualisere sig
almindelig i den hele Natur, og hvad betyder dette Ord uden at antage den
mest bestemte og uforanderlige Form? Denne Drivt synes \emph{Schelling} at
forvirre med Lyst til Afsondring, Isolation eller Uafhængighed af det
Absolute, med den Synken fra Ideernes ned i Sandsernes Verden, der skal
være den egentlige Kilde til al physisk Udartning og moralsk Fordervelse.
Men Ensligheden bestaaer ei i noget sandseligt og foranderligt; Grundformen
tilligemed de Regler, hvorefter den hos ethvert Individ udvikler sig, er en
evig Idee. Ved at uddanne sig efter dens Mönster forlader ikke Sielen sit
Udspring, men vender meget mere tilbage. Man undskylde, at jeg her bruger
Mysticismens Talemaader, hvilken jeg for Resten her hverken agter at
bestride eller forsvare, da det kun er min Hensigt at vise, at min Mening ei
heller strider mod denne Skoles Grundsætninger. At den med en mere empirisk
Philosophie stemmer overens, ja, at man, for at begribe hvad Erfarenhed
lærer, ei kan undvære den, er oven for allerede beviist.

% printed p. 247 (PDF 257)
Der gives altsaa et fast Begreb om enslige Ting. Dette Begreb har objectiv
Gyldighed eller Realitet. Det er altsaa med de menneskelige Siele, med
organiske Legemer, der kan betragtes som hines udvortes anskuelige Billeder,
langt anderledes beskaffen end med raae Stene eller Skyer og
Dunstsamlinger, om hvilke man kun i Almindelighed kan sige hvad de ere, men
ei anföre noget besynderligt Kiendetegn, som jo, medens de endnu ere til,
kan forandres til de modsatte. Imidlertid maae man tilstaae, at hine
Gienstande, ligesaavel som disse, tilsidst synes at forsvinde. Individuerne
forgaae medens Arten vedvarer. Lad dem end i det sidste Öieblik af deres
Tilværelse beholde samme væsentlige Form som i det förste, lad end Ideen
leve, saa er den virkelige Gienstand dog ikke mere. Theorien er altsaa,
kunde man slutte, om ei urigtig, i det mindste uden Nytte. Det er Umagen
værd at undersöge, om det forholder sig saa eller ikke.

De Ideer, som opstaae i den menneskelige Forstand, de Planer, den udkaster,
have ei altid nogen virkelig Gienstand. Tilligemed Tegningen eller Udkastet
er ikke Verket selv strax færdigt. Forstanden kan være frugtbar, men
Frugterne hvile, saa at sige, i dens eget Skiöd; uden for sig synes den saa
lidet at kunne frembringe nogen Ting, at ei engang Villie og Forsæt deraf
ere uudeblivelige Fölger. For at realisere sin Idee behöver den, som man i
Almindelighed troer, ei alene andre, men endog fremmede Kræfter, som
undertiden ere vanskelige nok at bevæge til en med den harmonisk
Medvirkning. Er Arbeidet vel begyndt, kan uforudseete Hindringer dog
afbryde, en fiendtlig Kraft nedbryde det. Mangler i Anlægget, som vi selv
op%
% printed p. 248 (PDF 258)
dage, kan giöre, at vi slaae Formen itu for efter en rigtigere Plan at
begynde paa en nye. De fleste Menneskeverker ere Korthuse, som vi vel i
Begyndelsens Hede kan have bestemt til evig Afmindelse, men som vi dog selv
bagefter indtil de sidste Spor tidt söge at udrydde. Men skulde det samme
vel ogsaa gielde om Naturens og Forsynets Verker, eller just det modsatte
deraf være rigtigt? Skulde her nogensinde mangle enten Midler eller Evne
til at udföre dem, eller mangen Bygning maaskee anlagt efter en umoden Plan
af en mere oplyst Forstand tidt blive forstyrret, for af dens Materialier at
opreise en fuldkomnere og til Hensigten mere passende? Enten man antager en
indvortes Nödvendighed eller en uendelig Visdom, som hvad den hele Verdens
Indretning angaaer har anlagt og udförer alting paa det bedste; saa maae
Ordenens Love sættes som uforanderlige, de stridende Kræfter omsider bringe
hverandre selv til et Slags svævende Ligevægt, Forvirring oplöse sig til
Harmonie, og deri maae den almindelige Tendents bestaae. Vel grunde disse
Sætninger sig paa eller före omsider til et intellectualt System i den hele
Natur, alligevel paatrænge Erfarenhed og Iagttagelser selv enhver Grandsker
saa uimodstaaelig en Overbeviisning om deres Rigtighed, at endog de
Philosopher, der ei have skammet sig ved at regne Lykken med blandt de
virkende Naturaarsager, dog have ladet den handle paa en gandske
regelmæssig Maade. Et stort Beviis derpaa ere blandt andre Democrits
utallige Verdener, hvori de samme individuelle Væsener bestandig forekomme.

Tingenes Grundformer ere fölgelig ligesom de platoniske Ideer
uforanderlige, ja, de eneste virkelige Ting, og disse
% printed p. 249 (PDF 259)
ere, som jeg tilforn har viist, ei Arter og Slægter som Gienstande for de
almindelige Begreber, i det mindste ei dem allene, men Individuer. At
Arternes Bestandighed altid har ligesom hört til de almindelige
Troesartikle, det har dog vel foruden Erfarenhed ogsaa i Naturlovenes
Uforanderlighed havt en höiere theoretisk Grund, uden hvilken en altid
ufuldstændig Induction dog ingen Vished kunde give. Denne Grund overföre vi
nu med Rette fra Arterne til Individuerne. Thi endskiöndt Erfarenhed er saa
langt fra her at stadfæste den, at de Ensliges Undergang meget mere hvert
Öieblik synes at vidne derimod; saa reiser sig dette Skin dog virkelig kun
af Misforstaaelse; hvilket tydelig kiendes, naar vi blive Principet troe.
Det er beviist, at ethvert Individ har en Grundform, hvis Hovedlinier under
alle dets afvexlende Tilstande, hvis vort Blæk var dybt nok, vilde befindes
uudslettelige. Disse ere nemlig kun synlige i den foranderlige Stof, som
beklæder og tillige fordunkler dem, ligesom de Træk, der skrives med det
sympathetiske Blik nogenledes blive læselige, naar en anden Materie faaer
givet dem Farve; men ere de vel mindre virkelige, fordi de af sig selv ei
falde i Öinene, eller fordi det ei kan skee uden at Formen derved lider
nogen Forandring?

De fleste Naturkyndige have lignet den organiske Stof, hvoraf Dyr og
Planter i Moders Liv blive dannede, ved en Kime, eller tillagt den et Slags
Præformation, d.~e.\ en indvortes Beqvem%
% printed p. 250 (PDF 260)
hed til den bestemte Form, der maaskee efter Aartusindes Forlöb först
bliver istand til at udvikle sig. Den organiske Kraft er efter Stoffens
Beskaffenhed forskiellig i hvert Legeme, ei regellös virkende, men ifölge
Regler, der for ethvert Tilfælde paa det nöieste ere bestemte; thi
endskiöndt udvortes Aarsager, der stöde til, kan forandre dens Retning; saa
kan de dog ei saaledes tilintetgiöre dens Virksomhed, at dennes særegne
Tendents for et opmærksomt Öie jo bliver kiendelig. Ingen Udartning er saa
stor, intet Misfoster selv saa vanskabt, at den oprindelige Form aldeles er
forvendt. Det er besynderligt, at baade Stoikerne og andre Materialister
have antaget dette Præformations System, men ikke desmindre kan have troet,
at ved Döden bliver denne Form saaledes brudt eller knust, at en Fornyelse
deraf aldrig mere kan finde Sted. Kan den da efter det grovere Legemes
Oplösning ei ligesaa vel vedvare, som den för Födselen, uvist hvorlænge,
under saa mange og store Ödelæggelser usynlig har vidst at holde sig
vedlige? Thi de störste Omvæltninger i Naturen lade de egentlige Kimer selv
være ubeskadede, ja befordre selv undertiden deres Væxt: de fineste og os
ofte forborgne Oplösninger skaane dem uden Tvivl ligeledes eller finde i
disse mindste Legemer maaskee en Modstand, som svarer til deres Finhed, og
staaer i omvendt Forhold til deres Störrelse. Man vilde meget misforstaae
mig, om man troede, at jeg ved disse Kimer eller Individuernes Grundform
forstaaer andet end et Anlæg eller en
% printed p. 251 (PDF 261)
vis Bestemmelse i den dannende Kraft, der er forskiellig i ethvert Individ:
Dette troer jeg enhver Naturforsker, hvad enten han med Bonnet antager
Evolutions-Systemet, eller med de fleste nyere et Slags Epigenesis, meget
vel kan indrömme. I övrigt vilde det före mig for vidt, hvis jeg her
omstændeligere skulde vise hvorledes jeg forestiller mig denne Sag.
Stoikernes \emph{rationes seminales} vare ligeledes kun Ævner og Anlæg til
visse Former, ingenlunde disse Former selv i det Smaae.

Fra den Standpunkt, paa hvilken vi ved disse Undersögelser ere komne, kan
hvert Menneske betragtes som en af de mange muelige Former, hvilke deres
fælles Natur eller Begreb kan modtage. Personer ere fölgelig for det förste
ei ubetydelige Smaadele af et Heelt, men enhver er dette Hele selv under sin
særegne Form. Man bör dernæst ei ansee Individuerne for mange mere eller
mindre fuldkomne Aftryk af samme Original, hvoraf erhvert maatte være saa
meget mindre værd som Mængden er större. Thi, naar af saadan Mængde kun en
eneste Copie blev tilbage, er der i en vis Henseende ved de övriges
Undergang kun lidet tabt. Men udgiöre derimod Individuer ikke blot et Tal,
hvori alle Enheder ere lige; saa kan deres Vedligeholdelse, hvad deres
væsentlige Form angaaer, ei være mindre vigtig end Arternes og Slægternes.
Vi maae fölgelig formode, at intet Individ virkelig forgaaer, men at Döden i
Gierningen kun er én Forvandling eller Forberedelse dertil; hvis Begyndelse
og Frem%
% printed p. 252 (PDF 262)
gang det dog ligesaa lidt er os tilladt at forfölge med vore Iagttagelser,
som hvad der för Undfangelsen med Kimen foregaaer.

Betragter man Naturen kun som en Samling af Kunststykker opstillede til
Beskuelse, fordrer dog Fuldstændighed, at intet maae mangle. Da Naturen
aldrig tillader sig noget Spring, men Rekken af dens Producter overalt maae
være uafbrudt; saa kan ligesaa lidt noget Individ, som nogen Slægt, Art
eller Afart være borte. Men Naturen er noget mere end blot et saadant
Galerie: den er et fuldkommen organisk Heelt, for hvilket enhver af de
selvstændige Dele, der udgiöre det, have samme Nödvendighed eller Vigtighed,
som det Hele selv eller alle övrige for denne. Endog den mindstes Tab kan ei
erstattes ved nogen anden, fordi netop dennes bestemte Form allene passer
ind i det Sted, hvor den staaer. Man kan herved hverken sammenligne
Menneskers ufuldkomne Verker eller engang Naturens egne, saaledes som de i
Tiden eller for Öieblikket i Erfarenhed fremstille sig. I begge hænder det
vel, at en Deel uden synderligt Tab kan ombyttes med en anden, ja, at Verket
derved endog kan blive forbedret; men Ideer ere, ligesom Kunstens Idealer,
uforgængelige: og endskiöndt man, i det man stræber at naae dem, stedse
retter og forandrer, saa gives der dog baade i Henseende til Delene og det
Hele noget characteristisk, der saavel for ethvert Stykke som for Arten
afgiver en Regel, fra hvilken Naturen selv aldrig tör afvige.

% printed p. 253 (PDF 263)
\emph{Aristoteles} var, som bekiendt, af den Mening, at kun det Almindelige
og Store kan være Gienstanden for det guddommelige Forsyn; at dette efter
visse Love regierer det Hele; men at det Besynderlige, de enslige Ting enten
ere alt for ubetydelige til at være dette Forsyn værdige, eller dets store
Plan saa hinderlige, at den uden deres Opofrelse ei kan udföres. Uden just
at höre til Peripatetikernes Skole have mange andre baade i Theorien antaget
samme Grundsætning, og endnu tiere i Handlemaade rettet sig derefter. I
Statskunsten har den stedse været det store Hiul, der har bestemt Maskinens
Gang. Men, da de, der styre den, ansee sig selv som Artens Repræsentanter;
saa finde de deri Grund til en dem fordeelagtig Undtagelse. Denne Grund er
alligevel gandske falsk, da de Selskaber hine forestille, som
\emph{Pomponatius} i hans Bog om Besværgelser allerede har bemærket,
ingenlunde udgiöre Arten selv, men ere ligesaa forgængelige som enslige
Personer, og, saavidt de ere menneskelige Indretninger, ingen fast Charakter
eller sand Individualitet kan have.

Man kunde troe, at den modsatte Mening, som jeg her har sögt at forsvare,
förer til en i Praxis endnu mere skadelig moralsk Egoisme. Denne Fölge vilde
være gandske rigtig, dersom enten Individuet nogensinde kunde forgaae, eller
det ved en uegennyttig Stræben for andres Fuldkommenhed blev hindret i sin
egen Uddannelse. Men saaledes forholder det sig ikke. Ideer%
% printed p. 254 (PDF 264) — final page of the essay
nes Gienstande ere, ligesom de selv, evige: og af Sædelighed, som er den
almindelige Fornuftslov, der byder indvortes og udvortes Harmonie, have
begge al deres Skiönhed og Glands.

Dog, Anvendelsen af denne Lærdom i Omgang med vore Lige, i det offentlige og
private Liv behövede jeg i et populært Foredrag maaskee udförlig at tilföie,
endskiöndt den ei er vanskelig at giöre. Men i et lærd Selskab, hvor det
lettere bör forbigaaes og de mest indviklede Knuder, saa vidt mueligt er,
alene blive löste, er det nok at pege paa Anvendelsen. En mere pragmatisk
Udförelse, der gaaer ind i det Besynderlige og Concrete, er derfor hverken
nödvendig eller paa dette Sted engang gandske passende.

\begin{center}\rule{0.25\linewidth}{0.4pt}\end{center}

\end{document}
